\scanpage{44}
On \deleted{\unsure{procède} de façon \unsure{analogue}} peut se ramener, grâce au th de structure, à prouver que une extension d'un schéma abélien $A$ par $G_a$ n'est pas rigide [\unsure{prouver|rem.} de \unsure{suite} qu'un quotient d'un rigide est rigide]. \uline{Or} \uline{ça n'est pas vrai} : une extension \uline{non triviale} de $A$ par $G_a$, du moins en car. 0, est rigide.

\uline{Proposition}. Soit $G$ schéma en groupes commutatif \inserted{loc} de type fini sur $S$, \uline{rigide}. Soit $n$ un entier premier aux car. des corps résiduels de $S$. Alors \deleted{les noyaux $G$} pour tt ouvert $U$ de $G$, $U$ est le plus petit ss-schéma fermé de $U$ qui majore les $ {}_{n^k}G\cap U$.

Démonstration donnée plus haut.

\uline{Proposition} $G$ un $S$-préschéma, en gr. \unsure{non} comm., \deleted{\unsure{plat} \unsure{et}}\inserted{\unsure{séparé}}, de type fini sur $S$, \deleted{$S$ \unsure{noeth.}} Soit $n$ un entier \deleted{premier} premier aux car des corps résiduels de $S$. Alors \inserted{le morphisme $G\xrightarrow{n}G$} \inserted{est \deleted{\unsure{étale}}\inserted{non ramifié}, en particulier le} schéma $ {}_nG$ est \deleted{étale}\inserted{non ramifié} sur $S$ ; \deleted{de \unsure{m.}} si $G$ est plat sur $S$, $ {}_nG$ est $=$ étale sur $S$.
% archive sheet number: 43
